% Track d6 -- Duplication in motion: pair creation and the balance surface.
% Verification: investigations/rest_energy/verify/d6_pair_balance.py (claims d6-0 ... d6-14).
% New references are in sections/references.tex.
\section{Duplication in motion: the balanced pair}
\label{d6:sec:main}

Two further remarks are tested here: that the $b$-particle of \cref{d4:sec:bfield} becomes decisive once two or more manifolds are
present, and that one should find the process by which a unitary (undivided) energy results in a duplication, from $\Cl(3,3)$ to
$\Cl(6,6)$. Duplication is taken literally: two copies of the arena whose tangent spaces add, $V_2=V\oplus V\cong\R^{6,6}$, which is grade
one of $\Cl(6,6)\cong\Cl(3,3)\hat\otimes\Cl(3,3)$ (equation~(2) of \cite{ES17}; the graded tensor product is inherited). It is followed in
motion, inside the closure of Section~17 of \cite{ES17}, and at a single field point. In brief:
(1) with one constituent and Lorentz-class inertia a constant nonzero rate forces a null momentum field, while with two or more a
codimension-two balance surface carries massive constituents at an exactly constant rate with bounded frames (\cref{d6:thm:balance},
restating Proposition~17.18 of \cite{ES17} with the dimension count added). This is the precise sense in which a duplicate does what one copy
cannot, and it needs no $b$-field. (2) The closure never creates a pair (\cref{d6:cor:nopair}). (3) The admissible $b$-field between the
manifolds, the copy rotation, duplicates a field literally, and no Lorentz-covariant operation makes a massive pair out of a single null
field (\cref{d6:prop:covariant}); a natural candidate that rotates a single direction between the copies is a different, non-covariant
element (\cref{d6:ex:R1}).
(4) The rest energy of a duplicated field point is $M_{\rm pair}=\half(\sigma_1+\sigma_2)$, and the copy rotation can move all of it into the two
copies (\cref{d6:thm:pair,d6:ex:bw}): the kinematics of $\gamma\gamma\to e^+e^-$, inherited. (5) The rest energy of each copy separately is
not invariant under the copy rotation, hence gauge dependent if that rotation is gauged, and with equal inertias the closure's own symmetry maps every balanced pair to one null constituent
(\cref{d6:prop:gauge}): pair creation is physical only where the copies are distinguished. (6) On the Dirac side the graded duplication adds
Dirac masses in quadrature (\cref{d6:prop:dirac}). (7) The zero-point energy stays zero, but an inter-manifold $b$ at unequal rates can
bring one line of the motion to zero exactly, and there the motion is unbounded (\cref{d6:prop:zpe}).

\paragraph{Conventions.} As in \cref{d3:sec:main} for one constituent. For two, $V_2=V\oplus V$ with coordinates $(X^{(1)};X^{(2)})$,
$\eta_2=\id_2\otimes\eta$, $G_2=\id_2\otimes G$; in the tower $e^{(1)}_a=e_a\otimes\id$, $e^{(2)}_a=\omega\otimes e_a$, $\omega=e_1\cdots e_6$,
and $V_2$ is spanned by the twelve $e^{(k)}_a$. $Z=(z_1,z_2)\in\C^2\otimes\C^3$, $z_k=z(X^{(k)})$, so $G_2$ acts as $i$; per copy
$E_6^{(k)},E_0^{(k)},\mathbf S^{(k)}$ and $q_k=z_k\cdot z_k$ are those of \eqref{d1:eq:defs}, and $q_{\rm tot}=q_1+q_2$. The Gram matrix
$Q_{kl}=z_k\cdot z_l$ (complex bilinear, symmetric) has singular values $\sigma_1\ge\sigma_2\ge0$; $H_{\rm sum}=\sum_kz_kz_k^\dagger$. The
\emph{copy rotation} is generated by $E:=\epsilon\otimes\id_6$, $\epsilon=\bigl(\begin{smallmatrix}0&-1\\1&0\end{smallmatrix}\bigr)$:
$R_\varphi=e^{\varphi E}$ maps $(z_1,z_2)\mapsto(\cos\varphi\,z_1-\sin\varphi\,z_2,\ \sin\varphi\,z_1+\cos\varphi\,z_2)$, identically in all six
directions; $E^2=-\id$, $[E,G_2]=0$, $E$ is antisymmetric for $\eta_2$ and for the Euclidean form. $T(z)$ is the stress tensor of
\cref{d1:prop:rainich}, $P(t)=T(z)t$ the density seen by a unit future timelike observer $t$, $P_{\rm sum}(t)=\sum_kT(z_k)t$, and
\[
M_{\rm pair}:=\inf_t|P_{\rm sum}(t)|,
\]
the rest energy of the duplicated field point. In the closure the per-copy quantities are those of the momentum fields
$\pi^{(k)}=\eta^{-1}p^{(k)}$, $w_k=q(z(\pi^{(k)}))$; inertias are of Lorentz class, $N_k=(\nu_{k,1}\eta^{-1}+\nu_{k,2}G\eta^{-1})/m_k$, with
$c_k:=\nu_k/m_k\in\C$ and $R_{\rm eff}e^{i\varphi_{\rm eff}}:=\sum_kc_kw_k(0)$ (Proposition~17.18 of \cite{ES17}). Two Lorentz actions occur: the
per-copy algebras $\lG^{(1)}\oplus\lG^{(2)}$, and the diagonal $\mathfrak l_{G,\Delta}$ ($k_B=2$ in Theorem~10.8 of \cite{ES17}). Each $E_0^{(k)}$ is invariant under both; $M_{\rm pair}$, defined with one observer for both copies, is invariant under the diagonal action but not under the per-copy one, since a boost of one copy alone changes it. A constituent is called \emph{massive} when its momentum field is not null, $E_0(\pi^{(k)})\neq0$:
massive in the sense of the dictionary of \cref{s0:sec:question}, an interpretive reading that does not involve $mc^2=\hbar\kappa$. Rates are rates:
$mc^2=\hbar\kappa$ is an added identification (\cref{d3:sec:main}). Identifiers d6-$n$ refer to \texttt{verify/d6\_pair\_balance.py}
(integrations DOP853, relative tolerance $10^{-12}$; d6-0 checks the conventions).

\subsection{Mass with a constant rate needs a second constituent}
\label{d6:sec:balance}

\begin{theorem}[Constant rate with massive constituents]
\label{d6:thm:balance}
Let one rotor act on $n$ constituents with Lorentz-class inertias and all $c_k\neq0$.
\begin{enumerate}[label=(\roman*),leftmargin=2.2em]
\item Along every solution $w_k(\tau)=e^{2i\theta(\tau)}w_k(0)$, so $R_{\rm eff}$ is constant. The rate $h/I_0$ is constant and nonzero iff
$R_{\rm eff}=0$ and $h(0)\neq0$; then $\theta=\kappa\tau$ and every frame moves on the circle
$x_k(\tau)=x_k(0)+\bigl(e^{2\kappa\tau G}-\id\bigr)(2\kappa G)^{-1}N_kp^{(k)}(0)$, of diameter $|N_kp^{(k)}(0)|/|\kappa|$.
\item For $n=1$, $R_{\rm eff}=0$ iff $E_0(\pi)=0$: no datum with a constant nonzero rate has a massive momentum field.
\item For $n\ge2$ the balance set $\mathcal B=\{\sum_kc_kw_k=0\}\subset\R^{6n}$ is, away from $0$, a smooth submanifold of codimension two; its
part on which every $E_0(\pi^{(k)})$ vanishes has dimension $4n<6n-2$, so the data with a constant nonzero rate and some $E_0(\pi^{(k)})>0$ are
open and dense in $\mathcal B\times\{h\neq0\}$. For $n=2$ both constituents are then massive, with $|c_1|E_0(\pi^{(1)})=|c_2|E_0(\pi^{(2)})$.
\end{enumerate}
\end{theorem}

\begin{proof}
(i) With a common rotor the momenta evolve as $p^{(k)}=e^{-\theta G}p^{(k)}(0)$ (Proposition~17.8 and the closing remark of Theorem~17.7 of
\cite{ES17}, $\mathfrak s_n=\R\ltimes_{G_n}V_n$), so $w_k=e^{2i\theta}w_k(0)$ (\cref{d3:thm:frozen}) and, by Theorem~17.11(ii) of \cite{ES17} per
constituent, $U(\theta)=\sum_k\half p^{(k)\top}N_kp^{(k)}=\half\Re\bigl(e^{2i\theta}\sum_kc_kw_k(0)\bigr)$, whence
$I_0\ddot\theta=R_{\rm eff}\sin(2\theta+\varphi_{\rm eff})$. If $R_{\rm eff}=0$, $\ddot\theta=0$; if $R_{\rm eff}>0$ and $h$ is constant, then
$\sin(2\theta+\varphi_{\rm eff})\equiv0$, so $\theta$ is constant and $h=0$. Lorentz-class $N_k$ anticommutes with $G$, so
$\dot x_k=e^{\theta G}N_ke^{-\theta G}p^{(k)}(0)=e^{2\theta G}N_kp^{(k)}(0)$ (Theorem~17.12 of \cite{ES17}), which integrates to the circle.
(ii) $R_{\rm eff}=|c|\,|w(p)|=2|c|E_0(\pi)$. (iii) $F(p)=\sum_kc_kw(p^{(k)})$ has $dF_p(\delta)=2\sum_kc_k\,z(\pi^{(k)})\cdot z(\delta\pi^{(k)})$. At
$p\neq0$ some $z(\pi^{(k)})\neq0$, and $u\mapsto2c_kz(\pi^{(k)})\cdot u$ maps $\C^3$ onto $\C$, so $0$ is a regular value of $F$ on
$\R^{6n}\setminus0$. The null set $\{w=0\}\subset\R^6$ is the cone $|\bB|=|\bE|$, $\bB\cdot\bE=0$, of dimension four, and the all-null part of
$\mathcal B$ is its $n$-fold product. Points with $w(p^{(1)})=1$, $w(p^{(2)})=-c_1/c_2$ show that the massive part is not empty.
\end{proof}

Items (i) and (ii) restate Proposition~17.18 and Corollary~17.14 of \cite{ES17} (fixed there); (iii) is fixed, the dimension count being added.
Checks: $\half p^{\!\top}Np=\Re(\nu w(p))/(2m)$ exactly, and the rank of $dF$ at random balanced points for $n=1,2,3$ (d6-1a,b); integrations of
the full two-constituent closure, rotor, momenta and frames (d6-2a--d). With equal inertias and $p^{(2)}=-Gp^{(1)}$ ($E_0(\pi^{(k)})=4.95$) the
rate stays constant to $1.6\times10^{-14}$ over $\tau\in[0,60]$ and the frames lie on the circles of (i) to $4\times10^{-12}$; a generic balanced
pair, and unequal inertias ($m=1,2.5$; $\nu=1+0.3i,\,0.4-0.7i$), give the same to $10^{-13}$. Off the surface ($R_{\rm eff}=0.0752$, rotating) the
rotor follows $I_0\ddot\theta=R_{\rm eff}\sin(2\theta+\varphi_{\rm eff})$ to $2\times10^{-10}$ and frame~1 drifts by $0.1795$ per half turn, exactly
additively (Theorem~17.16 of \cite{ES17} for several constituents).

\begin{remark}[The Lorentz requirement is what empties (ii)]
\label{d6:rem:lorentzreq}
$N=\operatorname{diag}(1,-1,0,-1,1,0)/m$ anticommutes with $G$ but is not $\lG$-invariant. For $p_0=e^1+e^2$, $z(\pi_0)=(1,1,0)$ and
$E_0(\pi_0)=1$, while $U(\theta)=\half p_0^{\!\top}e^{2\theta G}Np_0\equiv0$: the rate is exactly constant and the frame moves on a circle (d6-1c). One
constituent can thus combine a constant rate, a massive momentum field and a bounded, moving frame only through a cancellation between directions
that breaks Lorentz invariance: a bounded clamped frame needs $N_cp_0=0$, and a nonzero Lorentz-class $N_a$ then keeps $U$ constant only on null
data. The balance surface is the Lorentz-invariant form of the same cancellation, between constituents. Without the bounded frame no cancellation
is needed: a duality-class inertia, $[N,G]=0$, gives a constant rate for every datum, massive ones included, with $U=\half p_0^{\!\top}Np_0$ constant,
and the frame drifts at the velocity $Np_0$ (\cref{d3:cor:flat}(iii) and Theorem~17.12 of \cite{ES17}; restated). Fixed, as a counterexample to
(ii) without the hypothesis.
\end{remark}

\begin{corollary}[The closure creates no pair]
\label{d6:cor:nopair}
For every Hamiltonian on $\mathfrak s_n^*$, including any coupling of the constituents through $H$, $p^{(k)}(\tau)=e^{-\theta(\tau)G}p^{(k)}(0)$ with
one angle $\theta$. Hence $E_6$, $E_0$ and $\mathbf S$ of each $\pi^{(k)}$, the Hermitian Gram matrix $\bar z_k\cdot z_l$ and the moduli $|Q_{kl}|$
of the momentum fields are constants of motion, and $Q\mapsto e^{2i\theta}Q$. No closure Hamiltonian turns a null constituent into a massive
one, fills an empty copy, or moves rest energy between constituents.
\end{corollary}

\begin{proof}
The Lie--Poisson equations of $\mathfrak s_n$ are \eqref{d3:eq:LPgen} with $G_n=\bigoplus_kG$; the first integrates to $p=e^{-\theta G_n}p_0$,
$\dot\theta=H_h$.
\end{proof}

Obstructed, for the momentum field and every $H$; it extends \cref{d3:thm:frozen} to $n$ constituents. With a random Hamiltonian that couples
the two copies (quartic in $h$, trigonometric in $p$) these quantities stay constant to $2\times10^{-13}$ and $Q$ follows $e^{2i\theta}Q$ to the
same accuracy (d6-3). Inside the closure a pair is therefore a label of the initial data, carried and not generated, and the balance surface is
where such labels are compatible with a constant rate.

\subsection{The balanced pairs}
\label{d6:sec:pairs}

\begin{proposition}[Balanced pairs]
\label{d6:prop:pairs}
Let $z_k=z(\pi^{(k)})$.
\begin{enumerate}[label=(\roman*),leftmargin=2.2em]
\item Equal inertias. With $a=\half(z_1-iz_2)$ and $b=\half(z_1+iz_2)$, $q_1+q_2=4\,a\cdot b$ and $q_1=a\cdot a+2a\cdot b+b\cdot b$. The balance
surface is the quadric $a\cdot b=0$, and on it $q_2=-q_1$ and $E_0^{(1)}=E_0^{(2)}=\half|a\cdot a+b\cdot b|$.
\item $EG_2=\epsilon\otimes G$ is an $\eta_2$-antisymmetric involution, the noncompact generator of $\mathfrak g(0,2,0)\cong\so{2,\C}$. Its
eigenspaces $V_\pm=\{X^{(2)}=\pm GX^{(1)}\}$ are six-dimensional and totally $\eta_2$-null, and $E=\mp G_2$ on $V_\pm$. A momentum field lies in
$V_+$ iff $p^{(2)}=-Gp^{(1)}$ iff $b=0$, and in $V_-$ iff $p^{(2)}=Gp^{(1)}$ iff $a=0$; so $V_\pm$ lie in the balance surface, and on them the
copy rotation $R_\varphi$ is the duality rotation $e^{\mp\varphi G_2}$ of both copies.
\item On $V_+$, $z_2=iz_1$: $w_2=-w_1$, and the two momentum fields have equal $E_6$, $\mathbf S$ and $E_0$; when $q_1>0$ the first is of
magnetic and the second of electric type. $T(z_2)=T(z_1)$, so the summed stress tensor is $2T(z_1)$ and $M_{\rm pair}=2E_0^{(1)}$, although
$q_{\rm tot}=0$.
\item Unequal inertias. If $q_1\neq0$ the balanced pairs are $z_2=\rho\Lambda z_1$ with $\Lambda\in SO(3,\C)$ and $\rho^2=-c_1/c_2$, and
$|c_1|E_0^{(1)}=|c_2|E_0^{(2)}$.
\end{enumerate}
\end{proposition}

\begin{proof}
(i) Expand. (ii) $(\epsilon\otimes G)^2=\epsilon^2\otimes G^2=\id$ and $(\epsilon\otimes G)^\dagger=\epsilon^{\!\top}\otimes G^\dagger=-\epsilon\otimes G$.
The eigenvalue $+1$ means $-GX^{(2)}=X^{(1)}$, i.e.\ $X^{(2)}=GX^{(1)}$, and then
$\eta_2(X,Y)=\eta(X^{(1)},Y^{(1)})+\eta(GX^{(1)},GY^{(1)})=0$ since $G$ is an anti-isometry; on $V_+$,
$E(X^{(1)},X^{(2)})=(-GX^{(1)},X^{(1)})=-G_2(X^{(1)},X^{(2)})$. Finally $\eta^{-1}p^{(2)}=G\eta^{-1}p^{(1)}$ iff $p^{(2)}=-Gp^{(1)}$, as $\eta G=-G\eta$, and
$z_2=iz_1$ iff $b=0$. (iii) $T$ is duality invariant (\cref{d1:prop:rainich}), and \cref{d1:prop:rainich} gives $|2T(z_1)t|=2E_0^{(1)}$ for
every $t$. (iv) $c_1q_1+c_2\rho^2q(\Lambda z_1)=c_1(q_1-q_1)=0$; conversely $q(z_2/\rho)=q_1\neq0$, and $SO(3,\C)$ is transitive on each complex
sphere $\{q=q_1\}$ (\cref{d1:thm:orbits}).
\end{proof}

The first sentence of (iii) restates the last sentence of Proposition~17.18 of \cite{ES17} (fixed there); (i), (ii) and (iv) are fixed (that
$V_\pm$ are null, and are exactly where the copy rotation is a duality rotation, could have failed); ``magnetic and electric pair'' and
``duality-balanced'' are readings. Checked exactly (d6-4a--c) and numerically (d6-4d). So the balanced pair of Proposition~17.18 of \cite{ES17} is
a magnetic and an electric constituent a quarter duality turn apart, with equal rest energies, whose duality phases the rotor turns together.
Its total $q$ vanishes, but $q_{\rm tot}$ is not the rest energy of the pair: its summed density has rest energy $2E_0^{(1)}$ for every
observer. Thus $\half|q_{\rm tot}|$ is what the rotor sees ($R_{\rm eff}=|\nu|\,|q_{\rm tot}|/m$ for equal inertias) and $M_{\rm pair}\ge\half|q_{\rm tot}|$ is
what an observer of the two copies sees (\cref{d6:thm:pair}). The block-diagonal inertia is the per-copy Lorentz requirement: the invariant
symmetric forms on $V_2^*$ number four under $\lG^{(1)}\oplus\lG^{(2)}$, six under $\mathfrak l_{G,\Delta}$ alone (cross terms,
$U=\half\Re\tr(CQ(\pi))$, $C$ complex symmetric, balance $\tr(CQ(\pi))=0$, again of codimension two) and two under $\Cent(G_2)\cap\so{6,6}$, all of
Lorentz class (d6-4e; the $n=2$ analogue of Theorem~17.24 of \cite{ES17}).

\subsection{What a copy rotation can and cannot duplicate}
\label{d6:sec:covariant}

\begin{proposition}[Lorentz-covariant duplication]
\label{d6:prop:covariant}
\begin{enumerate}[label=(\roman*),leftmargin=2.2em]
\item The real linear maps of $V_2$ commuting with $\mathfrak l_{G,\Delta}$ are the $A\otimes\id_3$, $A\in M_2(\C)$, acting on $Z$ with $G_2=i$. Inside
$\so{6,6}$ they form $\operatorname{span}\{E,EG_2\}\cong\so{2,\C}$, which is $\mathfrak g(0,2,0)$ of Theorem~10.8 of \cite{ES17}; its compact part
is $\{R_\varphi\}$.
\item $A\otimes\id_3$ maps $Q\mapsto AQA^{\!\top}$ and $Z\mapsto ZA^{\!\top}$ ($Z$ the $3\times2$ matrix of columns $z_k$). Hence $R_\varphi$ preserves
$H_{\rm sum}$, so every summed density $P_{\rm sum}(t)$, and $q_{\rm tot}$, $\det Q$, $\sigma_1$, $\sigma_2$. The unitary $A$, the group $U(2)$ of
energy-preserving duality rotations of two free Maxwell fields (end of Section~13.5 of \cite{ES17}; \cite{GaillardZumino81}), preserve
$H_{\rm sum}$, $\sigma_1$ and $\sigma_2$ but not $q_{\rm tot}$, and $\uu2\cap\so{6,6}=\R E$.
\item Every map $F:\C^6\to\C^6$, linear or not, commuting with the diagonal $SO(3,\C)$ sends $(z,0)$ with $q(z)=0\neq z$ to $(\alpha z,\beta z)$:
two null constituents. In particular $R_\varphi(z,0)=(\cos\varphi\,z,\sin\varphi\,z)$.
\end{enumerate}
\end{proposition}

\begin{proof}
(i) $V$ is an irreducible $\lG$-module with commutant $\operatorname{span}\{\id,G\}\cong\C$ (Proposition~17.4 of \cite{ES17}), so the commutant
of $\mathfrak l_{G,\Delta}$ on $\R^2\otimes V$ is $M_2(\R)\otimes\operatorname{span}\{\id,G\}=M_2(\C)$; antisymmetry for $\eta_2$ leaves $A\in\C\epsilon$, which
is Proposition~13.14 of \cite{ES17} (row ``diagonal $\lG$'', dimension $2$). (ii) $Q=Z^{\!\top}Z$ and $H_{\rm sum}=ZZ^\dagger$, and $T$ is linear in
$H$ (\cref{d1:prop:lagrange}). (iii) For $\Lambda\in SO(3,\C)$ fixing $z$, $F(z,0)=F(\Lambda z,0)=\Lambda F(z,0)$. In the action $u\mapsto w\times u$
of $\so{3,\C}$ (\eqref{d4:eq:w}) the stabiliser algebra of a null $z$ is $\{w:w\times z=0\}=\C z$, whose common fixed vectors satisfy
$z\times v=0$, i.e.\ $v\in\C z$.
\end{proof}

Item (i) restates Proposition~13.14 and Theorem~10.8 of \cite{ES17} (fixed there), (ii) is inherited linear algebra and restates the remark
on $U(n)$ at the end of Section~13.5 of \cite{ES17}, and (iii) is an obstruction; all checked exactly (d6-5a--c). This settles in the negative the reading of the copy rotation
as a pair-creating operation. The copy rotation acting identically in all six directions is the admissible $b$-field between the two manifolds,
and it duplicates a field literally: the undivided energy and the rest energy of $(z,0)$ are shared in the same ratio
$\cos^2\varphi:\sin^2\varphi$, a null field becomes two null fields and a massive one two massive replicas. It cannot produce a balanced massive
pair from a single null field, and no other Lorentz-covariant operation can.

\begin{theorem}[The rest energy of a duplicated field point]
\label{d6:thm:pair}
Let $z_1,z_2\in\C^3$, with Gram matrix $Q$ of singular values $\sigma_1\ge\sigma_2$.
\begin{enumerate}[label=(\roman*),leftmargin=2.2em]
\item $T(z_1)\cdot T(z_2)=|z_1\cdot z_2|^2$, and as mixed tensors $\bigl(T(z_1)+T(z_2)\bigr)^2=\tfrac14\|Q\|_F^2\,\id+T(z_1\times z_2)$, with
$(z_1\times z_2)\cdot(z_1\times z_2)=\det Q$.
\item $|P_{\rm sum}(t)|^2=\tfrac14\|Q\|_F^2+u_t(z_1\times z_2)$ for every observer $t$, $u_t(w)$ being the energy density of $w$ seen by $t$. Hence
$M_{\rm pair}=\half(\sigma_1+\sigma_2)$, attained exactly in the rest frames of $z_1\times z_2$ when it is not null.
\item In every frame, $E_0^{(1)}+E_0^{(2)}\le M_{\rm pair}\le|P_{\rm sum}(t)|\le E_6^{(1)}+E_6^{(2)}$, and $M_{\rm pair}\ge\half|q_{\rm tot}|$.
\item $\max_\varphi\bigl(E_0^{(1)}+E_0^{(2)}\bigr)(R_\varphi Z)=M_{\rm pair}$.
\item For two null fields $M_{\rm pair}=|z_1\cdot z_2|=\bigl(E_6^{(1)}E_6^{(2)}\bigr)^{1/2}(1-\cos\vartheta_{12})$, $\vartheta_{12}$ the angle between
their Poynting vectors; for $(z,0)$ with $z$ null, $M_{\rm pair}=0$.
\end{enumerate}
\end{theorem}

\begin{proof}
(i) Polynomial identities, verified symbolically in twelve real variables (d6-6a); in spinor form \cite{PenroseRindler84}, $T=\sum_k\phi_k\bar\phi_k$ and
$\phi^k_{AB}\phi^{l\,B}{}_C=\half(\phi^k\cdot\phi^l)\epsilon_{AC}+\chi^{kl}_{AC}$ with $\chi^{kl}$ antisymmetric in $(k,l)$, so the cross terms of $T^2$
pair the symmetric $Q$ with the antisymmetric $\chi$ and drop out. The last identity is the complex Lagrange identity. (ii) $T$ is
$g$-self-adjoint, so $|P(t)|^2=t\cdot T^2t$; insert (i). By \cref{d1:thm:min}, $u_t(w)\ge E_0(w)=\half|\det Q|$, with equality iff $t$ is a
rest frame of $w$, and $\tfrac14\|Q\|_F^2+\half|\det Q|=\tfrac14(\sigma_1+\sigma_2)^2$. For null $w\neq0$ the infimum is approached by Doppler boosts;
for $w=0$, $|P(t)|=\half\sigma_1$ for all $t$. (iii) Each $P_k(t)$ is future causal of length $E_0^{(k)}$ (\cref{d1:prop:rainich}), future causal
vectors obey $|P_1+P_2|\ge|P_1|+|P_2|$ and $|P|\le P^0$, and $\sigma_1+\sigma_2\ge|\tr Q|$. (iv) Write $q_{1,2}(\varphi)=t\pm d(\varphi)$ and
$Q_{12}(\varphi)=r(\varphi)$ with $t=\half q_{\rm tot}$. The pair $(d,r)$ turns by $2\varphi$, $d^2+r^2=t^2-\det Q$ is fixed, and
$(|q_1|+|q_2|)^2=2|t|^2+2|d|^2+2|t^2-d^2|\le2|t|^2+2|d|^2+2|r|^2+2|\det Q|=(\sigma_1+\sigma_2)^2$, by $|t^2-d^2|=|\det Q+r^2|\le|\det Q|+|r|^2$.
Equality needs $r(\varphi)^2$ to be a non-negative multiple of $\det Q$; as $\varphi$ varies, $r(\varphi)$ runs over an ellipse centred at $0$ or a
segment through $0$, so such $\varphi$ exists. (v) $Q=\bigl(\begin{smallmatrix}0&r\\r&0\end{smallmatrix}\bigr)$, $r=z_1\cdot z_2$; and
$T(z_k)=k_k\otimes k_k$ with $k_k$ future null, $k_k^0=(E_6^{(k)})^{1/2}$, so $|r|=k_1\cdot k_2$ by (i), and $|\mathbf S^{(k)}|=E_6^{(k)}$.
\end{proof}

Items (i), (iii) and (v) are inherited: ordinary electrodynamics of two field strengths, with the reverse triangle inequality as the kinematics
of invariant mass. The closed form of $M_{\rm pair}$ in (ii) and its saturation by the copy rotation in (iv) are fixed and new; (i), (ii) and
(iv) are not stated in \cite{ES17}. Checked symbolically (i), at the rest frames of $z_1\times z_2$ to $10^{-10}$ against $300$ random
observers (ii), by maximising over $\varphi$ (relative deviation $5\times10^{-16}$) (iv), and on random and null pairs (iii), (v) (d6-6a--d).
Since $T(z_1\times z_2)$ is traceless, the trace of (i) gives the pair invariant of \cref{d5:rem:twoE0},
$E_{\rm pair}=(\tfrac14\tr T_{\rm tot}^2)^{1/2}=\half\|Q\|_F$, and (ii) gives $M_{\rm pair}^2=E_{\rm pair}^2+\half|\det Q|$. The two agree exactly when
$\det Q=0$, which holds on every duplicate $(\alpha z,\beta z)$ of a single field (\cref{d5:thm:dup}); for the counter-propagating null pair of
\cref{d6:ex:bw}, $E_{\rm pair}=\sqrt2$ while $M_{\rm pair}=2$. The invariant mass of the pair is $M_{\rm pair}$, not $E_{\rm pair}$ (inherited).

\begin{example}[Two counter-propagating null fields]
\label{d6:ex:bw}
$z_1=(1,i,0)$ and $z_2=(1,-i,0)$ are null with $\mathbf S^{(1)}=-e_3=-\mathbf S^{(2)}$, so $\mathbf S_{\rm sum}=0$, $E_6^{\rm tot}=2$,
$Q=\bigl(\begin{smallmatrix}0&2\\2&0\end{smallmatrix}\bigr)$ and $M_{\rm pair}=2$. The rotation $R_{\pi/4}$ gives $z_1'=(0,\sqrt2\,i,0)$ and
$z_2'=(\sqrt2,0,0)$: a pure electric and a pure magnetic field, both at rest, with $q=-2,+2$ and $E_0=1$ each. In this frame, the centre-of-momentum frame, the chain of
\cref{d6:thm:pair}(iii) holds with equality for the rotated pair, $E_0^{(1)}+E_0^{(2)}=M_{\rm pair}=|P_{\rm sum}|=E_6^{\rm tot}=2$, while
$M_{\rm pair}=2>\half|q_{\rm tot}|=0$: the balanced pair is produced at rest in the centre-of-momentum frame and takes the whole invariant mass
(d6-7). Inherited: the kinematics of $\gamma\gamma\to e^+e^-$ \cite{BreitWheeler34}; the same superposition is the standing wave of
\cref{d1:prop:whole}(3).
\end{example}

\begin{example}[The rotation of one direction]
\label{d6:ex:R1}
Let $R^{(1)}_\varphi$ rotate the planes $(e^{(1)}_1,e^{(2)}_1)$ and $(e^{(1)}_4,e^{(2)}_4)$ through $\varphi$. It is compact, commutes with $G_2$ and lies in
$\Cent(G_2)\cap\so{6,6}\cong\so{6,\C}$, the $n=2$ analogue of Theorem~17.19 of \cite{ES17}. On $((1,i,0),0)$ it gives $z_1'=(\cos\varphi,i,0)$,
$z_2'=(\sin\varphi,0,0)$: $q_1'=-\sin^2\varphi$, $q_2'=\sin^2\varphi$, a balanced pair with $E_0=\half\sin^2\varphi$ each. $E_6^{\rm tot}=1$ and $q_{\rm tot}=0$
are conserved, $\mathbf S_{\rm sum}$ goes from $-e_3$ to $-\cos\varphi\,e_3$ and $M_{\rm pair}$ from $0$ to $\sin\varphi$, so
$E_6^{\rm tot}=1\ge M_{\rm pair}=\sin\varphi\ge\sum_kE_0^{(k)}=\sin^2\varphi$, with equality throughout at $\varphi=\pi/2$ (a magnetic field along $e_1$ and
an electric one along $e_2$, at rest). $R^{(1)}$ does not commute with $\mathfrak l_{G,\Delta}$, and lies neither in $\mathfrak g(0,2,0)$ nor in
$\mathfrak g(0,1,1)=\so{3,3}$ (Lorentz algebra on copy~1, copy~2 a spectator): it is a candidate gauge transformation under no Lorentz action of
Theorem~10.8 of \cite{ES17} with $k_I=0$ (d6-8), and under none with $k_I>0$ either: it lies outside $\mathfrak g(2,0,0)$ and
$\mathfrak g(1,0,1)=\so{5,3}$ (x6-8a), and it has components from the interior four-plane of one copy into the binding frame of the other,
which the proof of Theorem~10.8 excludes from $\mathfrak g(1,1,0)$. It is therefore non-covariant under all five Lorentz actions at $n=2$. Read kinematically it is an infinitely heavy third body, absorbing momentum and no energy, and
$E_6\ge E_0^{(1)}+E_0^{(2)}$ is the field-point form of the threshold $E_\gamma\ge2mc^2$ for pair production in the field of a nucleus
\cite{BetheHeitler34}. The numbers and the non-covariance are fixed (d6-8; x6-8a,b), the third-body reading is interpretive, the threshold
kinematics is inherited, and the agent is constructed: it exists, and nothing in the arena supplies it.
\end{example}

A tempting reading of pair creation applies to $((1,i,0),0)$ a rotation said to mix the copies ``identically in all six directions''. The rotation
that produces the balanced pair above is $R^{(1)}_\varphi$, which rotates two directions only, not the copy rotation $R_\varphi$. For $R^{(1)}_\varphi$ the
conservation of $E_6^{\rm tot}$ and $q_{\rm tot}$ and the non-conservation of $\mathbf S_{\rm sum}$ hold, while the copy rotation, which does act identically
in all six directions, conserves $\mathbf S_{\rm sum}$ and the whole summed stress tensor.

\subsection{Is the rest energy of one copy observable?}
\label{d6:sec:gauge}

\begin{proposition}[Gauge status of the per-copy rest energy]
\label{d6:prop:gauge}
\begin{enumerate}[label=(\roman*),leftmargin=2.2em]
\item $R_\varphi$ is admissible for the diagonal action: compact, commuting with $G_2$ and $\mathfrak l_{G,\Delta}$, and the compact part of
$\mathfrak g(0,2,0)$. On $S_2$ it is generated by $M=\half\sum_a\eta_{aa}e^{(1)}_ae^{(2)}_a$, with $\ad M=-E$ on $V_2$ and eigenvalues
$i\{-3,\dots,3\}$; $U_\varphi=\prod_a\bigl(\cos\tfrac\varphi2+\sin\tfrac\varphi2\,\eta_{aa}e^{(1)}_ae^{(2)}_a\bigr)$ has $U_{2\pi}=\id$, so the copy rotation
acts through $SO(2)$; and $U_{\pi/2}J^{(1)}U_{\pi/2}^{-1}=J^{(2)}$, $J^{(k)}=e^{(k)}_2e^{(k)}_3$.
\item $R_\varphi$ preserves $H_{\rm sum}$, $q_{\rm tot}$, $\det Q$, $\sigma_{1,2}$, $M_{\rm pair}$ and the total charge $J^{(1)}+J^{(2)}$; it does not
preserve $E_0^{(k)}$, $q_k$ or the per-copy charges (\cref{d6:ex:bw}, (i)). On $V_\pm$ it is a duality rotation of both copies, and there the
$E_0^{(k)}$ are invariant.
\item Equal inertias. The two-constituent closure is invariant under $O(\eta_2)\cap\Cent(G_2)\cong O(6,\C)$, whose compact part, the real
orthogonal group of $\C^6$, acts transitively on $\{q_{\rm tot}=0,\ |Z|=r\}$: every balanced pair is carried to one null constituent and an empty
copy at the same $E_6^{\rm tot}$; for instance $(e_1,ie_1)\mapsto((1,i,0),0)$, with $M_{\rm pair}$ going from $1$ to $0$.
\item If $c_1\neq c_2$, no nonzero element of $\mathfrak g(0,2,0)$ preserves the inertia.
\end{enumerate}
\end{proposition}

\begin{proof}
(i) Each term pairs two generators of the same square, hence is compact (Theorem~2.3 of \cite{ES17}), squares to $-\tfrac14$, and the six terms
commute; exponentiate term by term. (ii) \Cref{d6:prop:covariant}(ii) and \cref{d6:prop:pairs}(ii). (iii) Every $\Lambda$ commuting with $G_2$
preserves the Lie--Poisson structure, and for $\Lambda\in O(\eta_2)$, $\Lambda\eta_2^{-1}\Lambda^{\!\top}=\eta_2^{-1}$ and
$\Lambda G_2\eta_2^{-1}\Lambda^{\!\top}=G_2\eta_2^{-1}$, so the equal Lorentz-class inertia is invariant. A nonzero $Z$ with $q_{\rm tot}(Z)=0$ is $u+iv$
with $u,v\in\R^6$, $|u|=|v|$, $u\perp v$, and $SO(6)$ is transitive on such pairs. (iv) $sE+tEG_2$ acts on the copy index as $(s+it)\epsilon$, and
$\epsilon^{\!\top}\operatorname{diag}(c_1,c_2)+\operatorname{diag}(c_1,c_2)\,\epsilon=(c_2-c_1)\bigl(\begin{smallmatrix}0&1\\1&0\end{smallmatrix}\bigr)$, which must vanish for the form $\sum_kc_k\,z_k\cdot z_k$ to be invariant.
\end{proof}

Item (i) restates Theorem~10.8 and Proposition~10.9 of \cite{ES17} (fixed there), with the lift computed; (ii)--(iv) are fixed; checked in integer
arithmetic (d6-9, d6-10). This decides whether ``mass creation in pairs'' is physical. \emph{If the copy rotation is gauged}, the rest energy of
each copy is not an observable: \cref{d6:ex:bw} is a change of gauge, not a process, and two counter-propagating null fields and an
electric--magnetic pair at rest are one physical state, whose gauge-invariant content $M_{\rm pair}=2$ was there before. \emph{In the
equal-inertia closure alone}, its own compact symmetry maps every balanced pair to one null constituent, changing $M_{\rm pair}$ as it does so (by (iii)), and the rotor sees only
$\half|q_{\rm tot}|=0$; the mass of the pair is invisible to the rate. $M_{\rm pair}$ becomes an invariant only through the diagonal Lorentz action, and the per-copy rest energies only through a per-copy one,
neither of which the closure selects (the paper leaves the Lorentz action at several constituents open, Table~7 of \cite{ES17}; at $n=3$ the closure symmetry meets the interior Lorentz algebra of the first constituent in $0$, Theorem~17.20 there, and \cref{d5:prop:so6C}(iv) is the $n=2$ analogue). \emph{If the copies are
distinguished} --- by unequal inertias, which by (iv) break $\mathfrak g(0,2,0)$, by unequal binding rates, by a projector, or by a Lorentz
action on one copy only, whose algebra $\mathfrak g(0,1,1)$ contains no copy mixing --- the per-copy rest energies are observables, but then the
closure creates no pair (\cref{d6:cor:nopair}). In the interior with $S=\operatorname{diag}(\kappa_1,\kappa_2)\otimes G+\beta E$ the momentum field
obeys $\dot\pi=S^\dagger\pi$, i.e.\ $z(\pi)(\tau)=e^{ih\tau}z(\pi)(0)$ with $h=\bigl(\begin{smallmatrix}\kappa_1&-i\beta\\i\beta&\kappa_2\end{smallmatrix}\bigr)$ on the
copy index. For $\kappa_1=\kappa_2$ this is $e^{i\kappa\tau}R_{-\beta\tau}$, and the per-copy rest energies change only by the gauge rotation
(Larmor \cite{Larmor1897}; \cref{d4:sec:bfield}). For $\kappa_1\neq\kappa_2$ they change physically: for the pair of \cref{d6:ex:bw},
$E_0(\pi^{(1)})(\tau)=|z_1\cdot z_2|\,|U_{11}U_{12}|$, $U=e^{ih\tau}$, which for $\kappa=(1,3)$, $\beta=0.4$ peaks at $0.6897$, the Rabi value
$2\sqrt{s(1-s)}$ with $s=4\beta^2/((\kappa_2-\kappa_1)^2+4\beta^2)$, at fixed $E_6^{\rm tot}$; from $(z,0)$ with $z$ null both copies stay null for all
$\tau$ (d6-11). This is two-level mixing \cite{Rabi37}, inherited. It changes the per-copy rest energies but not $M_{\rm pair}$, so their sum stays at most
$M_{\rm pair}$ (\cref{d6:thm:pair}(iii)) without being conserved in general: for the pair of \cref{d6:ex:bw} both copies gain the same rest
energy, so none passes from one copy to the other, and from two null fields rest energy appears only where the second is non-collinear and
supplies a nonzero $M_{\rm pair}$.

\begin{remark}[A copy rotor]
\label{d6:rem:copyrotor}
The uniqueness argument of Theorem~17.7 of \cite{ES17} uses that equivariant maps between distinct constituents vanish, which holds for
$\lG^{(1)}\oplus\lG^{(2)}$ but not for $\mathfrak l_{G,\Delta}$: $E$ is $\mathfrak l_{G,\Delta}$-equivariant, commutes with $G_2$ and rotates $p$ without dilating it, so a
second rotor $\hat c$ with $[\hat c,X]=EX$, $[\hat c,\hat g]=0$ gives an equivariant closure $\R^2\ltimes V_2$, with Lie--Poisson equations
$\dot p=(H_hG_2^{\!\top}+H_{h_c}E^{\!\top})p$, $\dot h=-p^{\!\top}G_2\nabla H$, $\dot h_c=-p^{\!\top}E\nabla H$. With equal Lorentz-class inertias
$\dot h_c=0$, the copy angle turns uniformly, and the per-copy $E_0(\pi^{(k)})$ oscillate while the summed momentum tensor, $M_{\rm pair}(\pi)$ and
$|q_{\rm tot}(\pi)|$ stay constant; $((1,i,0),0)$ stays null--null (d6-12). The inter-manifold $b$-field made dynamical changes
the rest energies of the two manifolds at fixed pair invariants. Their sum is not conserved in general (for the data of \cref{d6:ex:bw} the copy
rotation takes it from $0$ to $2$) and never exceeds $M_{\rm pair}$ (\cref{d6:thm:pair}(iii)), so the rotor creates per-copy rest energy but no
pair invariant. Constructed.
\end{remark}

\subsection{The Dirac side}
\label{d6:sec:dirac}

\begin{proposition}[Pairs of Dirac states in the tower]
\label{d6:prop:dirac}
On $S_2=\R^8\otimes\R^8$ with the tower generators:
\begin{enumerate}[label=(\roman*),leftmargin=2.2em]
\item $J^{(1)}=J\otimes\id$ and $J^{(2)}=\id\otimes J$ (the Klein factors cancel on bivectors); they commute with each other and with the eight
interior gammas $e^{(k)}_a$, $a\in\{1,4,5,6\}$. With an external $i$, $Q_{\rm tot}=-i(J^{(1)}+J^{(2)})$ has spectrum $\{-2,0,2\}$ with
multiplicities $16,32,16$.
\item Let $\psi,\phi\in S\otimes\C$ be particles at rest, $e_1\psi=\psi$, $J\psi=i\psi$ ($H(0)=me_1$, equation~(8) of \cite{ES17}), and $C=e_2$
(Theorem~5.1 of \cite{ES17}). Then $\Psi=\psi\otimes C\phi$ has charges $(1,-1)$, $Q_{\rm tot}\Psi=0$, and
$-i\bigl(JH(0)\otimes\id+\id\otimes JH(0)\bigr)\Psi=2m\Psi$.
\item $C^{(1)}C^{(2)}=e^{(1)}_2e^{(2)}_2$ anticommutes with $J^{(1)}$ and $J^{(2)}$, but $(C^{(1)}C^{(2)})^2=-\id$: no state is mapped to $\pm$ itself.
The Klein-corrected $C\otimes C=-\omega^{(1)}C^{(1)}C^{(2)}$, $\omega^{(1)}=\omega\otimes\id$, squares to $\id$, and $\Psi_\pm=\psi\otimes C\psi\pm C\psi\otimes\psi$
satisfy $(C\otimes C)\Psi_\pm=\pm\Psi_\pm$ and $Q_{\rm tot}\Psi_\pm=0$.
\item The tower's mass terms anticommute: $(e^{(1)}_1+e^{(2)}_1)^2=2\,\id$, so $m(e^{(1)}_1+e^{(2)}_1)$ has spectrum $\{\pm\sqrt2\,m\}$, whereas the
two-particle operator $m(e_1\otimes\id+\id\otimes e_1)=m(e^{(1)}_1+\omega^{(1)}e^{(2)}_1)$ has spectrum $\{-2m,0,2m\}$, multiplicities $16,32,16$.
\item Under the diagonal interior action ($k_I=2$) a rotation through $2\pi$ is $+\id$ on $S_2$.
\end{enumerate}
\end{proposition}

\begin{proof}
(i) $(\omega\otimes e_2)(\omega\otimes e_3)=\id\otimes J$ as $\omega^2=\id$; the commutations are the blade rule (equation~(3) of \cite{ES17});
$(J^{(1)}+J^{(2)})^2=-2(\id-J^{(1)}J^{(2)})$, where $\half(\id-J^{(1)}J^{(2)})$ is a projector of trace $32$. (ii) $C$ anticommutes with $e_1$ and $J$, and
$J$ commutes with $H(0)$. (iii) Generators of distinct constituents anticommute and $e_2^2=\id$; $\id\otimes e_2=\omega^{(1)}e^{(2)}_2$. (iv) The
same anticommutation; $e_1\otimes\id$ and $\id\otimes e_1$ are commuting traceless involutions with traceless product. (v) Each copy's rotation
through $2\pi$ is $-\id$, and the two commute.
\end{proof}

Inherited: Klein--Jordan--Wigner signs \cite{JordanWigner28}, the charge-conjugation parity of a particle--antiparticle pair
\cite{WeinbergQFT}, and integer spin for two spin-$\half$ constituents, which is Proposition~10.9 of \cite{ES17} at $k_I=2$ (restated, fixed
there); checked exactly (d6-13a--e). Two points are specific to the tower. The total charge is invariant under the copy rotation while
$U_{\pi/2}$ exchanges $J^{(1)}$ and $J^{(2)}$ (\cref{d6:prop:gauge}(i)), so which copy carries the particle is gauge dependent exactly as the
per-copy rest energy is. And by (iv) the graded duplication $\Cl(3,3)\hat\otimes\Cl(3,3)=\Cl(6,6)$ adds Dirac masses in quadrature, like two
anticommuting Clifford directions --- one particle of mass $\sqrt2\,m$ --- while the pair threshold $2m$ appears only for Klein-corrected,
ungraded operators: duplication in the tower's sense is not pair creation.

\subsection{Zero-point energy of the duplicated frame}
\label{d6:sec:zpe}

\begin{proposition}[Zero-point energy]
\label{d6:prop:zpe}
\begin{enumerate}[label=(\roman*),leftmargin=2.2em]
\item For $\mathcal L=\sum_k\frac{m_k}2\eta\bigl(\dot X^{(k)}+\kappa_kGX^{(k)},\dot X^{(k)}+\kappa_kGX^{(k)}\bigr)$ each rate $\kappa_k$ carries Krein
signature $(3,3)$; with positive-norm quantisation the zero-point energy is $\sum_k\frac{\hbar\kappa_k}2(3-3)=0$ for all $\kappa_1,\kappa_2>0$ and $m_1,m_2\neq0$. At $\kappa_k=0$ the copy is free, its phase-space generator is not semisimple, and no zero-point energy is defined.
\item With an inter-manifold $b=\beta E$ ($m_1=m_2$, $\kappa_k\ge0$ not both zero), $S=\operatorname{diag}(\kappa_1,\kappa_2)\otimes G+\beta E$ has spectrum
$\pm i\mu_\pm$, $\mu_\pm=\half(\kappa_1+\kappa_2)\pm\bigl(\tfrac14(\kappa_1-\kappa_2)^2+\beta^2\bigr)^{1/2}$, each three times. Every solution is
bounded unless $\beta^2=\kappa_1\kappa_2$ and $\kappa_1\neq\kappa_2$; when bounded, each line has Krein signature $(3,3)$ and the zero-point energy
vanishes. At $\kappa_1=0$, $|\mu_-|\approx\beta^2/\kappa_2$.
\item At $\beta^2=\kappa_1\kappa_2$, $\kappa_1\neq\kappa_2$, $\mu_-=0$ and the phase-space generator is not semisimple: some solutions grow linearly
in $\tau$.
\end{enumerate}
\end{proposition}

\begin{proof}
(i) \Cref{d2:prop:framevac} per copy, or \cref{d4:thm:krein}. (ii) On the $G_2=i$ sector $S$ acts on the copy index as
$i\operatorname{diag}(\kappa_1,\kappa_2)+\beta\epsilon=i(K-i\beta\epsilon)$ with $K-i\beta\epsilon$ Hermitian. $S^\dagger=K\otimes G-\beta E$ commutes with
$G_2$, so for $\mu\neq0$ its eigenspaces lie in the $G_2=\pm i$ sectors, which are totally isotropic for $h$ (a vector with $GY=iY$ has
$\bE=-i\bB$); \cref{d4:thm:criterion}(ii) then gives boundedness except at $\mu_-=0$, and \cref{d4:thm:krein} the signatures. (iii) At
$\beta^2=\kappa_1\kappa_2$ the kernel of $S^\dagger$ consists of the $(v_1c,v_2c)$, $c\in\C^3$, with $v_2=-i\kappa_1v_1/\beta$, on which
$\eta_2=\Re\bigl((v_1^2+v_2^2)\,c\cdot c\bigr)$ and $v_1^2+v_2^2=v_1^2(1-\kappa_1/\kappa_2)\neq0$; so the kernel is not $\eta_2$-null and
\cref{d4:thm:criterion}(i) fails.
\end{proof}

Item (i) restates Theorem~18.6 of \cite{ES17} (fixed there) for each copy: duplication preserves the neutral signature, and the cancellation
survives unequal rates and masses. Since $E=C_v$, (ii) and (iii) specialise \cref{d5:thm:bfield2} to $m_1=m_2$. The lines in (ii) are inherited two-level mixing and restate
\cref{d5:thm:bfield2}(ii); the boundedness criterion restates \cref{d5:thm:bfield2}(v); the seesaw value $0.0495097568$ at
$(\kappa_1,\kappa_2,\beta)=(0,5,\half)$ against $\beta^2/\kappa_2=0.05$ restates \cref{d5:thm:bfield2}(vi). What is new here is the Krein signature
$(3,3)$ of each line and the vanishing zero-point energy at $n=2$, fixed by \cref{d4:thm:krein}. (iii) restates the Jordan block in the proof
of \cref{d5:thm:bfield2}(v) (fixed there); it is the inter-manifold analogue of the exceptional ratio of \cref{d4:thm:screen}. Checked numerically for the signatures (d6-14a,b) and with rational
data for semisimplicity: not semisimple at $(1,4,2)$, semisimple at $(1,4,1)$, $(2,2,2)$ and $(1,3,\tfrac25)$ (d6-14c). Read with the added
identification $mc^2=\hbar\kappa$ applied to the lines: an inter-manifold $b$ can give an unbound manifold the line $\beta^2/\kappa_2$ borrowed from the
binding of the other, or bring a line to zero exactly at the price of unbounded motion; it never changes a rate, the spectrum of the
self-adjoint part of the shift (\cref{d7:thm:schr}(iii)), and never produces a line without an existing binding, since at
$\kappa_1=\kappa_2=0$ the lines $\pm\beta$ are removed by $X\mapsto e^{-\beta E\tau}X$.

\subsection{Assessment}
\label{d6:sec:assess}

\begin{table}[t]
\centering
\small
\begin{tabular}{@{}>{\raggedright\arraybackslash}p{2.7cm}>{\raggedright\arraybackslash}p{7.2cm}>{\raggedright\arraybackslash}p{3.2cm}>{\raggedright\arraybackslash}p{1.4cm}@{}}
\toprule
claim & statement & class & check\\
\midrule
Thm.~\ref{d6:thm:balance}(i) & rate constant iff $R_{\rm eff}=0$; frames on circles & fixed (restated, Prop.~17.18) & d6-2a--d\\
Thm.~\ref{d6:thm:balance}(ii) & $n=1$: constant rate forces $E_0(\pi)=0$ & fixed (restated, Cor.~17.14) & d6-1a\\
Thm.~\ref{d6:thm:balance}(iii) & $n\ge2$: codimension-two balance set, massive data dense & fixed & d6-1b\\
Rem.~\ref{d6:rem:lorentzreq} & without the Lorentz requirement $n=1$ suffices & fixed (counterexample) & d6-1c\\
Cor.~\ref{d6:cor:nopair} & no closure Hamiltonian creates or moves a pair & obstructed & d6-3\\
Prop.~\ref{d6:prop:pairs} & $a\cdot b=0$; $V_\pm$ null, $E=\mp G_2$; magnetic--electric pair, $M_{\rm pair}=2E_0$, $q_{\rm tot}=0$ & fixed (part restated); readings interpretive & d6-4a--e\\
Prop.~\ref{d6:prop:covariant}(i,ii) & commutant $M_2(\C)$, $\mathfrak g(0,2,0)$, invariants of $R_\varphi$ and $U(2)$ & restated / inherited & d6-5a,b\\
Prop.~\ref{d6:prop:covariant}(iii) & no covariant map makes a pair from one null field & obstructed & d6-5c\\
Thm.~\ref{d6:thm:pair} & $M_{\rm pair}=\half(\sigma_1+\sigma_2)$, the chain, $\max_\varphi$ & (i), (iii), (v) inherited; (ii), (iv) fixed (new) & d6-6a--d\\
Ex.~\ref{d6:ex:bw} & two null fields $\to$ electric and magnetic at rest & inherited & d6-7\\
Ex.~\ref{d6:ex:R1} & rotation of one direction: balanced pair from a null field, covariant under none of the five Lorentz actions at $n=2$; third body & fixed (numbers, non-covariance); agent constructed; reading interpretive & d6-8; x6-8a,b\\
Prop.~\ref{d6:prop:gauge} & $R_\varphi$ admissible; per-copy $E_0$ not invariant; $O(6,\C)$ maps pairs to null fields & fixed (part restated) & d6-9, 10\\
interior & distinguished copies: per-copy rest energy changed by two-level mixing, sum $\le M_{\rm pair}$ & inherited & d6-11\\
Rem.~\ref{d6:rem:copyrotor} & a dynamical copy rotor changes $E_0^{(k)}$ and their sum at fixed pair invariants & constructed & d6-12\\
Prop.~\ref{d6:prop:dirac} & Klein signs, $C$-parity, masses in quadrature & inherited & d6-13a--e\\
Prop.~\ref{d6:prop:zpe} & zero-point energy $0$; exceptional ratio $\beta^2=\kappa_1\kappa_2$ & fixed ((i) restated; lines, boundedness and (iii) restate \cref{d5:thm:bfield2}) & d6-14a--c\\
answers & duplication as copy rotation, pair creation and balance & interpretive & ---\\
\bottomrule
\end{tabular}
\caption{The claims of \cref{d6:sec:main}, their classes and the checks of \texttt{d6\_pair\_balance.py} (d6-0 verifies the conventions; x6-8a,b are printed by \texttt{x6\_check.py}).}
\label{d6:tab:claims}
\end{table}

\emph{On the $b$-particle with two manifolds.} The remark is partly borne out. With two manifolds there is an admissible $b$-field between
them, the copy rotation, the compact part of $\mathfrak g(0,2,0)$; it commutes with the binding root $G_2$, and with the binding $K\otimes G$ only at equal rates (\cref{d5:thm:bfield2}(iii)), keeps every orbit bounded except when
$\beta^2=\kappa_1\kappa_2$ with $\kappa_1\neq\kappa_2$, or when $\kappa_1=\kappa_2=0$ and nothing binds (\cref{d5:thm:bfield2}(v)), and changes the rest energy of each manifold, their sum bounded by $M_{\rm pair}$ but not conserved in general, which on one manifold only the anti-diagonal part of $b$
can do, and only by tearing or by screening (\cref{d1:prop:bfield,d4:thm:screen}); with unequal rates it mixes them, lending a line to an unbound
manifold or bringing one to zero, while the rates themselves never change (\cref{d7:thm:schr}(iii)). It does not generate an invariant mass: the copy rotation preserves $M_{\rm pair}$, the summed stress tensor and $q_{\rm tot}$, and the
interior flow $e^{ih\tau}$, unitary on the copy index, preserves $M_{\rm pair}$, the summed stress tensor and $E_6^{\rm tot}$ at any rates but
$|q_{\rm tot}|$ only at equal rates (for the pair of \cref{d6:ex:bw} at $\kappa=(1,3)$, $\beta=0.4$ it leaves $0$); with equal rates it is removable; the rest energy of each copy is not invariant under it, hence gauge dependent if it is gauged; and the one thing two constituents do in
the closure that one cannot, carrying rest energy at a constant rate, needs no $b$-field at all: it is the balance surface.

\emph{On how a unitary energy results in a duplication.} The mathematical process is the graded tensor product (inherited), under which the
undivided energy becomes $E_6^{\rm tot}=\half|Z|^2$; the Lorentz-covariant, energy--momentum-preserving duplications form $U(2)\otimes\id_3$, of
which the metric ones form $O(2)\otimes\id_3$ (the real orthogonal $A$); its identity component, the copy rotation $SO(2)$, is the part generated by an admissible $b$-field, and the other component contains the swap of the copies and $\operatorname{diag}(1,-1)$, which are metric but not copy rotations (\cref{d6:prop:covariant}). The physical process it
describes shares the energy of one copy between two as literal replicas, preserving the type. From one null field, a balanced magnetic--electric pair requires either
a second, non-collinear null field, with the inherited kinematics of $\gamma\gamma\to e^+e^-$ ($\sum_kE_0^{(k)}\le M_{\rm pair}$, saturated), or a
non-covariant agent that takes the momentum, with the kinematics of pair production on a nucleus ($\sum_kE_0^{(k)}\le E_6^{\rm tot}$). From one massive field $z$ the copy rotation never balances $(z,0)$, since it preserves $q_{\rm tot}=q(z)\neq0$; the non-metric equal split with a relative quarter turn of duality, $z\mapsto(z,iz)/\sqrt2$ (\cref{d5:thm:dup}(ii) at $\theta=\pi/2$), is covariant and energy-preserving, lands on $V_+$ and gives the balanced pair of \cref{d6:prop:pairs}(iii), with $E_0^{(k)}=\half E_0(z)$ each and $M_{\rm pair}=E_0(z)$ unchanged (\cref{d5:prop:balance}). In motion
the closure conserves the content of each copy and carries the pair, on the balance surface, at a constant rate. The decisive caveat is
\cref{d6:prop:gauge}: every statement that a mass has been created in one copy needs a structure that distinguishes the copies, and the
symmetric closure and the gauged copy rotation provide none. Whether such a structure is physical is open.
